\chapter{Converting phenotype-regression moments to breeding-value parameters}
\label{conversion:chapter}

A published selection regression usually describes relative fitness as a
function of measured phenotype. The variance share in this project instead
uses log expected output conditional on breeding value. Their coefficients
cannot generally be equated. Under stated conditions, the phenotype
regression supplies selected moments, those moments can be transferred to
breeding-value coordinates, and a log-quadratic model can be chosen to
reproduce them. This chapter derives each operation and identifies when that
model also recovers the actual conditional log-fitness derivatives.

The calculations use the tower property from
Chapter~\ref{conditioning:chapter}, Gaussian fourth moments from
Chapter~\ref{gaussian:chapter}, and the Gaussian selection formulas in
Chapter~\ref{selection:chapter}. Matrix order matters throughout: covariance
and curvature matrices need not commute.

\section{Define the regression and its reference population}
\label{conversion:regression}

Let \(\vect Z\) be the measured phenotype, let
\(\vect x=\vect Z-\E[\vect Z]\), and write
\(\mat P=\Cov(\vect Z)\). Let \(W\geq0\) be output in the same population and
selection episode, with \(0<\bar W=\E[W]<\infty\), and define relative
fitness \(w=W/\bar W\). Thus \(\E[w]=1\).

Consider the population least-squares regression
\begin{equation}
 w=\alpha_w+\vect\beta^{\mathsf T}\vect x
     +\frac12\vect x^{\mathsf T}\bgamma\vect x+\varepsilon,
 \qquad \bgamma=\bgamma^{\mathsf T}.
 \label{conversion:regression_model}
\end{equation}
We use \(\alpha_w\) to distinguish this relative-fitness intercept from the
log-mean intercept in Chapter~\ref{estimation:chapter}. This is the
linear--quadratic phenotype regression of selection analysis
\citep{LandeArnold1983}. Assume \(\E[w^2]<\infty\), so that the population
least-squares problem and the Gaussian-weighted moments used below are
well defined. The normal equations for least squares give
\begin{equation}
 \E[\varepsilon]=0,\qquad
 \E[x_i\varepsilon]=0,\qquad
 \E[x_ix_j\varepsilon]=0
 \quad\hbox{for every }i,j.
 \label{conversion:normal_equations}
\end{equation}
These equations require orthogonality to the fitted regressors. They do not
require \(\E[\varepsilon\mid\vect x]=0\) at every phenotype. Consequently
the moment calculations below can hold when the fitted polynomial is not the
exact conditional mean.

The factor \(1/2\) applies to the entire quadratic form. Expanding it gives
\begin{equation}
 \frac12\vect x^{\mathsf T}\bgamma\vect x
 =\frac12\sum_j\gamma_{jj}x_j^2
   +\sum_{j<k}\gamma_{jk}x_jx_k.
 \label{conversion:coefficient_convention}
\end{equation}
Thus a software coefficient \(c_{jj}\) multiplying \(x_j^2\) corresponds
to \(\gamma_{jj}=2c_{jj}\). A coefficient multiplying the single
cross-product column \(x_jx_k\) corresponds to \(\gamma_{jk}\) itself:
the two symmetric matrix entries have already combined. This convention
must be checked before using published coefficients
\citep{StinchcombeEtAl2008}. The coefficients must also use compatible trait
units and the same population centring.

\section{Recover selected phenotype moments}
\label{conversion:phenotype_moments}

Assume now that \(\vect x\sim N(\vect0,\mat P)\). For this section, the
formulas do not require inverses, although an identifiable full quadratic
regression is most conveniently written on the positive-variance support.

\begin{proposition}[Moments supplied by the phenotype regression]
\label{conversion:regression_moments}
Under the regression orthogonality and Gaussian assumptions,
\begin{align}
 \alpha_w&=1-\frac12\tr(\bgamma\mat P),\nonumber\\
 \E[\vect xw]&=\mat P\vect\beta,\nonumber\\
 \E[\vect x\vect x^{\mathsf T}w]
   &=\mat P+\mat P\bgamma\mat P.
 \label{conversion:phenotype_moment_formula}
\end{align}
The selected phenotype covariance is consequently
\begin{equation}
 \mat P_{\rm sel}
 =\mat P+\mat P
       (\bgamma-\vect\beta\vect\beta^{\mathsf T})\mat P.
 \label{conversion:selected_phenotype_covariance}
\end{equation}
\end{proposition}

\begin{proof}
First take expectations in
Equation~\eqref{conversion:regression_model}. The linear term and residual
have mean zero, while
\[
 \E[\vect x^{\mathsf T}\bgamma\vect x]
 =\sum_{j,k}\gamma_{jk}\E[x_jx_k]
 =\sum_{j,k}\gamma_{jk}P_{jk}
 =\tr(\bgamma\mat P).
\]
Since \(\E[w]=1\), this proves the intercept formula. It is a normalisation of
the fitted average over the population, not a statement that fitted relative
fitness at the mean phenotype equals one.

Next multiply the regression by \(x_i\) and average:
\begin{align*}
 \E[x_iw]
 &=\alpha_w\E[x_i]
   +\sum_j\beta_j\E[x_ix_j]
   +\frac12\sum_{j,k}\gamma_{jk}\E[x_ix_jx_k]
   +\E[x_i\varepsilon]\\
 &=\sum_jP_{ij}\beta_j.
\end{align*}
The mean and all third centred Gaussian moments vanish, and the residual
term vanishes by the normal equations. In vector form this is
\(\E[\vect xw]=\mat P\vect\beta\).

For the second moment, multiply by \(x_ix_l\). The linear and residual
terms again vanish, leaving
\[
 \E[x_ix_lw]
 =\alpha_wP_{il}
   +\frac12\sum_{j,k}\gamma_{jk}\E[x_ix_lx_jx_k].
\]
The fourth-moment identity proved in
Equation~\eqref{gaussian:fourth_moment} is
\[
 \E[x_ix_lx_jx_k]
 =P_{il}P_{jk}+P_{ij}P_{lk}+P_{ik}P_{lj}.
\]
Evaluate its three contributions separately:
\begin{align*}
 \frac12\sum_{j,k}\gamma_{jk}P_{il}P_{jk}
 &=\frac12P_{il}\tr(\bgamma\mat P),\\
 \frac12\sum_{j,k}\gamma_{jk}P_{ij}P_{lk}
 &=\frac12(\mat P\bgamma\mat P)_{il},\\
 \frac12\sum_{j,k}\gamma_{jk}P_{ik}P_{lj}
 &=\frac12(\mat P\bgamma^{\mathsf T}\mat P)_{il}
  =\frac12(\mat P\bgamma\mat P)_{il}.
\end{align*}
Symmetry of \(\mat P\) and \(\bgamma\) is used in the last two lines.
Thus
\[
 \E[x_ix_lw]
 =\left\{\alpha_w+\frac12\tr(\bgamma\mat P)\right\}P_{il}
   +(\mat P\bgamma\mat P)_{il}
 =P_{il}+(\mat P\bgamma\mat P)_{il}.
\]
Finally, relative-fitness weighting already has normalisation one.
Its covariance is its second moment minus the outer product of its mean:
\begin{align*}
 \mat P_{\rm sel}
 &=\mat P+\mat P\bgamma\mat P
   -(\mat P\vect\beta)(\mat P\vect\beta)^{\mathsf T}\\
 &=\mat P+\mat P
       (\bgamma-\vect\beta\vect\beta^{\mathsf T})\mat P.
\end{align*}
\end{proof}

These are moments under the nonnegative observed relative-fitness weights
\(w\). The fitted polynomial is used to calculate those moments through
orthogonality. It need not be positive everywhere and should not be
substituted as a probability weight over all phenotypes merely because its
regression coefficients are available.

\section{Condition breeding value on phenotype}
\label{conversion:conditioning}

Let \(\vect X=\vect A-\E[\vect A]\) be centred breeding value, with covariance
\(\mat G\). Assume \((\vect X,\vect x)\) is jointly Gaussian in compatible
trait coordinates, and
\begin{equation}
 \Cov(\vect X,\vect x)=\mat G.
 \label{conversion:cross_covariance}
\end{equation}
For now take \(\mat P\) positive definite. The cross-covariance assumption is
additional information about the trait decomposition; it does not follow
from fitting a phenotype regression.

\begin{proposition}[Conditional breeding-value moments]
\label{conversion:conditional_moments}
Under these assumptions,
\begin{align}
 \E[\vect X\mid\vect x]&=\mat G\mat P^{-1}\vect x,\nonumber\\
 \Cov(\vect X\mid\vect x)
 &=\mat G-\mat G\mat P^{-1}\mat G.
 \label{conversion:conditional_formula}
\end{align}
\end{proposition}

\begin{proof}
Define the regression residual
\[
 \vect U=\vect X-\mat G\mat P^{-1}\vect x.
\]
Its covariance with phenotype is
\[
 \Cov(\vect U,\vect x)
 =\mat G-\mat G\mat P^{-1}\mat P=\mat0.
\]
Because \(\vect U\) and \(\vect x\) are linear transformations of a joint
Gaussian vector, they are jointly Gaussian. Zero cross-covariance therefore
makes them independent. This follows, for example, by factoring the joint
Gaussian characteristic function with a block-diagonal covariance matrix.
Consequently \(\E[\vect U\mid\vect x]=\vect0\).

Expanding its covariance, with the matrix order retained, gives
\begin{align*}
 \Cov(\vect U)
 &=\mat G-\mat G\mat P^{-1}\mat G
       -\mat G\mat P^{-1}\mat G
       +\mat G\mat P^{-1}\mat P\mat P^{-1}\mat G\\
 &=\mat G-\mat G\mat P^{-1}\mat G.
\end{align*}
Since \(\vect X=\mat G\mat P^{-1}\vect x+\vect U\), conditioning fixes the
first term and leaves the independent residual covariance unchanged. This
proves both statements.
\end{proof}

The matrix \(\mat G\mat P^{-1}\) need not be symmetric. Its transpose is
\(\mat P^{-1}\mat G\), and these products must not be exchanged. For example,
the conditional second moment is
\begin{align}
 \E[\vect X\vect X^{\mathsf T}\mid\vect x]
 &=\mat S_c+\mat C\vect x\vect x^{\mathsf T}\mat C^{\mathsf T},\nonumber\\
 \mat S_c&=\mat G-\mat G\mat P^{-1}\mat G,\qquad
 \mat C=\mat G\mat P^{-1}.
 \label{conversion:conditional_second}
\end{align}
The two outer factors in the second term are \(\mat C\) and
\(\mat C^{\mathsf T}\), not two identical copies of \(\mat C\).

\section{Transfer reproductive weights to breeding-value moments}
\label{conversion:breeding_moments}

We now need a condition that connects the phenotype regression to the
reproduction of individuals with different breeding values. A sufficient
condition is
\begin{equation}
 \E[W\mid\vect X,\vect x]=\E[W\mid\vect x].
 \label{conversion:mediation}
\end{equation}
It says that, once the measured phenotype is given, breeding value supplies
no further information about expected reproductive output. Full conditional
independence is stronger than needed. For the following calculation, it
would also suffice to justify directly the corresponding conditional
first- and second-moment factorizations.

\begin{proposition}[Selected breeding-value moments]
\label{conversion:selected_breeding_moments}
Under the regression, Gaussian, cross-covariance and selection-mediation
conditions above, write
\(\mat K=\bgamma-\vect\beta\vect\beta^{\mathsf T}\).
Then
\begin{equation}
 \vect m=\E_{\rm sel}[\vect X]=\mat G\vect\beta,\qquad
 \Gsel=\mat G+\mat G\mat K\mat G.
 \label{conversion:selected_formula}
\end{equation}
\end{proposition}

\begin{proof}
The conditional mean condition implies
\begin{align*}
 \E[\vect Xw\mid\vect x]
 &=\E\{\vect X\,\E[w\mid\vect X,\vect x]\mid\vect x\}\\
 &=\E[\vect X\mid\vect x]\,\E[w\mid\vect x].
\end{align*}
Average across phenotypes and substitute
Equation~\eqref{conversion:conditional_formula}:
\begin{align*}
 \E_{\rm sel}[\vect X]
 &=\E[\vect Xw]\\
 &=\mat G\mat P^{-1}
        \E\{\vect x\,\E[w\mid\vect x]\}\\
 &=\mat G\mat P^{-1}\E[\vect xw]\\
 &=\mat G\mat P^{-1}\mat P\vect\beta
 =\mat G\vect\beta.
\end{align*}
The same mediation condition applies to
\(\vect X\vect X^{\mathsf T}\). Using
Equation~\eqref{conversion:conditional_second}, its selected second moment
is
\begin{align*}
 \E_{\rm sel}[\vect X\vect X^{\mathsf T}]
 &=\mat S_c\,\E[w]
   +\mat C\E[\vect x\vect x^{\mathsf T}w]\mat C^{\mathsf T}\\
 &=\mat G-\mat G\mat P^{-1}\mat G\\
 &\quad+\mat G\mat P^{-1}
       (\mat P+\mat P\bgamma\mat P)
       \mat P^{-1}\mat G\\
 &=\mat G-\mat G\mat P^{-1}\mat G
       +\mat G\mat P^{-1}\mat G+\mat G\bgamma\mat G\\
 &=\mat G+\mat G\bgamma\mat G.
\end{align*}
Subtract the outer product of the selected mean:
\begin{align*}
 \Gsel
 &=\mat G+\mat G\bgamma\mat G
      -(\mat G\vect\beta)(\mat G\vect\beta)^{\mathsf T}\\
 &=\mat G+\mat G
       (\bgamma-\vect\beta\vect\beta^{\mathsf T})\mat G.
\end{align*}
\end{proof}

\paragraph{Biological interpretation.}
Individuals with the same phenotype can have different breeding values.
Their conditional covariance remains part of the selected second moment.
Discarding it would treat the prediction \(\mat G\mat P^{-1}\vect x\) as an
observed breeding value and underestimate the unresolved variation.
Selection mediation is also substantive: if unmeasured traits, shared
environments or genetic effects predict reproduction after the measured
phenotype is given, the factorizations used above need not hold.

\section{Invert the selected moments}
\label{conversion:inversion}

Suppose \(\mat G\) and the selected covariance \(\Gsel\) are positive
definite. The exact Gaussian log-quadratic selection model in
Proposition~\ref{selection:selected_moments} has
\[
 \Gsel=(\mat G^{-1}-\mat H)^{-1},
 \qquad
 \vect m=\Gsel\vect b.
\]
These equations can be solved without assuming that their matrices commute.

\begin{proposition}[Moment-matched log-quadratic parameters]
\label{conversion:parameter_conversion}
Under the selected-moment conditions above and positive definiteness,
\begin{equation}
 \vect b=\Gsel^{-1}\mat G\vect\beta,\qquad
 \mat H=\mat G^{-1}-\Gsel^{-1},
 \qquad
 \Gsel=\mat G+\mat G
       (\bgamma-\vect\beta\vect\beta^{\mathsf T})\mat G.
 \label{conversion:conversion_formula}
\end{equation}
They define an integrable Gaussian log-quadratic weight with the specified
selected mean and covariance.
\end{proposition}

\begin{proof}
Invert the selected-covariance identity:
\[
 \Gsel^{-1}=\mat G^{-1}-\mat H.
\]
Rearranging gives \(\mat H=\mat G^{-1}-\Gsel^{-1}\).
Left-multiplying \(\vect m=\Gsel\vect b\) by \(\Gsel^{-1}\) gives
\(\vect b=\Gsel^{-1}\vect m=\Gsel^{-1}\mat G\vect\beta\).
Both inverse covariances are symmetric, so \(\mat H\) is symmetric.

To verify existence of the corresponding weight, multiply the original
normal density by \(\exp(\vect b^{\mathsf T}\vect X+
\vect X^{\mathsf T}\mat H\vect X/2)\). Its exponent, apart from a constant,
is
\begin{align*}
 -\frac12\vect X^{\mathsf T}(\mat G^{-1}-\mat H)\vect X
       +\vect b^{\mathsf T}\vect X
 &=-\frac12\vect X^{\mathsf T}\Gsel^{-1}\vect X
       +\vect m^{\mathsf T}\Gsel^{-1}\vect X\\
 &=-\frac12(\vect X-\vect m)^{\mathsf T}
          \Gsel^{-1}(\vect X-\vect m)
       +\frac12\vect m^{\mathsf T}\Gsel^{-1}\vect m.
\end{align*}
Positive definiteness of \(\Gsel\) makes this a finite Gaussian integral.
After normalisation, the weighted distribution is \(N(\vect m,\Gsel)\).
\end{proof}

The order in the slope formula is specifically
\(\Gsel^{-1}\mat G\vect\beta\), not
\(\mat G\Gsel^{-1}\vect\beta\). For another way to see the order, factor
\[
 \Gsel=\mat G(\mat I+\mat K\mat G)
      =(\mat I+\mat G\mat K)\mat G.
\]
The first factorisation gives
\[
 \Gsel^{-1}\mat G=(\mat I+\mat K\mat G)^{-1},
 \qquad
 \vect b=(\mat I+\mat K\mat G)^{-1}\vect\beta.
\]
Replacing \(\mat K\mat G\) by \(\mat G\mat K\) changes the expression unless
the matrices happen to commute. Neither commutation nor simultaneous
diagonalisation is assumed in the derivation.

Positive selected covariance is sufficient for an integrable matched model.
It does not require stabilising curvature. If \(\Gsel\preceq\mat G\), inverse
order gives \(\Gsel^{-1}\succeq\mat G^{-1}\), so \(\mat H\preceq0\);
other selected covariance changes can give other curvature signs.
Relative fitness cannot recover the absolute output multiplier:
adding a constant to the matched log mean leaves the selected moments
unchanged.

\subsection{Exact derivatives versus matched moments}

If the true \(W_A\) is exactly log quadratic in the breeding-value coordinates,
then its Gaussian selection formulas must have the same mean and covariance,
and the inversion recovers its \(\vect b,\mat H\). If exact log-quadraticity
fails while the moment conditions remain valid, the construction still gives
a log-quadratic weight with matching first and second selected moments.
Those integrated moments do not determine the local derivatives of an
arbitrary function \(\log\E[W\mid\vect A]\).

The distinction also applies to the variance share calculated from the
converted parameters. Under exact closure it describes the stated
log-quadratic relationship. Under moment matching it describes the chosen
quadratic model, whose agreement in selected moments does not establish
agreement in its complete distribution of log fitness.

\section{A positive fitness model that satisfies the assumptions}
\label{conversion:positive_model}

The supplied checks include positive fitness functions, so the conversion can
be assessed without treating a possibly negative regression polynomial as a
fitness function. Let centred phenotype be
\(\vect Z=\vect A+\vect E\), where
\(\vect A\sim N(\vect0,\mat G)\) and
\(\vect E\sim N(\vect0,\mat V_E)\) are independent. Take
\begin{equation}
 f_Z(\vect z)
 =\exp\{c+\vect q^{\mathsf T}\vect z
                  -\tfrac12\vect z^{\mathsf T}\mat C\vect z\},
 \qquad \mat C=\mat C^{\mathsf T}\succeq0,
 \label{conversion:positive_phenotype}
\end{equation}
with \(\mat V_E\) positive definite. This is a positive phenotype-dependent
conditional mean; reproduction may have additional stochastic variation
around it.

For a fixed breeding value \(\vect a\), write
\(\mat Q=\mat V_E^{-1}+\mat C\). In the environmental integral, collect
the terms involving \(\vect e\):
\begin{align*}
 \vect q^{\mathsf T}(\vect a+\vect e)
 -\frac12(\vect a+\vect e)^{\mathsf T}\mat C(\vect a+\vect e)
 -\frac12\vect e^{\mathsf T}\mat V_E^{-1}\vect e
 &=
 \vect q^{\mathsf T}\vect a-\frac12\vect a^{\mathsf T}\mat C\vect a\\
 &\quad-\frac12\vect e^{\mathsf T}\mat Q\vect e
       +(\vect q-\mat C\vect a)^{\mathsf T}\vect e.
\end{align*}
Completing the square in \(\vect e\) and using the Gaussian normalising
integral gives
\begin{align}
 \log W_A(\vect a)
 &=c-\frac12\log\det(\mat V_E)-\frac12\log\det(\mat Q)
     +\vect q^{\mathsf T}\vect a
     -\frac12\vect a^{\mathsf T}\mat C\vect a\nonumber\\
 &\quad+\frac12(\vect q-\mat C\vect a)^{\mathsf T}
                    \mat Q^{-1}(\vect q-\mat C\vect a).
 \label{conversion:multivariate_environment}
\end{align}
The positive-definite \(\mat Q\) ensures that the integral exists. Expand the
last term before reading off the coefficients:
\[
 \frac12\vect q^{\mathsf T}\mat Q^{-1}\vect q
 -\vect a^{\mathsf T}\mat C\mat Q^{-1}\vect q
 +\frac12\vect a^{\mathsf T}\mat C\mat Q^{-1}\mat C\vect a.
\]
Consequently,
\begin{equation}
 \vect b=\vect q-\mat C\mat Q^{-1}\vect q,\qquad
 \mat H=-\mat C+\mat C\mat Q^{-1}\mat C.
 \label{conversion:environment_parameters}
\end{equation}
These formulas provide an independent environmental-averaging route to the
same target. The reproduction code fits the phenotype regression by
quadrature and compares its moment conversion with this conditional
expectation in one and two traits. The finite comparison checks the
implementation. The agreement under the stated model follows from the
derivations, including all Gaussian and mediation assumptions.

\section{Centring, units and singular support}
\label{conversion:coordinates}

\subsection{Centre at the population mean}

The vanishing third moments used earlier refer to centred Gaussian
phenotypes. If a raw polynomial is initially expressed about a different
reference point, write its coordinate as \(\vect z=\vect x+\vect\delta\).
Then
\begin{align*}
 \alpha_0+\vect\beta_0^{\mathsf T}\vect z
       +\tfrac12\vect z^{\mathsf T}\bgamma\vect z
 &=\alpha_0+\vect\beta_0^{\mathsf T}\vect\delta
       +\tfrac12\vect\delta^{\mathsf T}\bgamma\vect\delta\\
 &\quad+(\vect\beta_0+\bgamma\vect\delta)^{\mathsf T}\vect x
       +\tfrac12\vect x^{\mathsf T}\bgamma\vect x.
\end{align*}
The centred linear coefficient is
\(\vect\beta=\vect\beta_0+\bgamma\vect\delta\); it is not generally
\(\vect\beta_0\). The regression space is unchanged by the shift when all
constant, linear and quadratic terms are retained, but its coefficients
change.

\subsection{Use compatible trait coordinates}

For an invertible coordinate transformation \(\vect x'=\mat T\vect x\),
the same polynomial has
\[
 \vect\beta'=\mat T^{-\mathsf T}\vect\beta,\qquad
 \bgamma'=\mat T^{-\mathsf T}\bgamma\mat T^{-1}.
\]
If breeding values are expressed in the same new coordinates,
\[
 \mat G'=\mat T\mat G\mat T^{\mathsf T},\qquad
 \mat P'=\mat T\mat P\mat T^{\mathsf T}.
\]
Substitution into the selected-moment formula gives
\(\vect m'=\mat T\vect m\) and
\(\Gsel'=\mat T\Gsel\mat T^{\mathsf T}\). Inverting the latter,
\[
 (\Gsel')^{-1}
 =\mat T^{-\mathsf T}\Gsel^{-1}\mat T^{-1},
\]
so the converted parameters transform as
\[
 \vect b'=\mat T^{-\mathsf T}\vect b,\qquad
 \mat H'=\mat T^{-\mathsf T}\mat H\mat T^{-1}.
\]
Thus the conversion is compatible with a coordinated change of trait units.
Using standardised regression coefficients with an untransformed
\(\mat G\) is not such a coordinated change.

\subsection{Restrict inverse operations to varying directions}
\label{conversion:support}

The moment formulas involving only multiplication remain meaningful when
covariances are singular, but inverses must be taken on the appropriate
support. First remove zero-variance phenotype directions. If
\(\Var(\vect v^{\mathsf T}\vect x)=0\), the centred phenotype component is
zero almost surely. The cross-covariance condition then gives
\(\vect v^{\mathsf T}\mat G\vect v=0\), so breeding value has no variation
in that same direction. An orthonormal basis for the phenotype support
therefore retains all relevant breeding-value variation and makes the
reduced \(\mat P\) invertible.

Breeding-value covariance can still have a smaller rank. Let \(\mat U\)
have orthonormal columns spanning the range of \(\mat G\). Then
\(\vect X=\mat U\vect X_s\) almost surely, and
\(\mat G_s=\mat U^{\mathsf T}\mat G\mat U\) is positive definite.
After calculating the selected moments, restrict them as
\[
 \vect m_s=\mat U^{\mathsf T}\vect m,\qquad
 \mat S_s=\mat U^{\mathsf T}\Gsel\mat U.
\]
If \(\mat S_s\) is positive definite, the identifiable parameters are
\[
 \vect b_s=\mat S_s^{-1}\vect m_s,\qquad
 \mat H_s=\mat G_s^{-1}-\mat S_s^{-1}.
\]
They describe the function on the breeding-value support. Derivatives in
directions with no breeding-value variation are not determined by these
moments. If selected covariance is degenerate within an initially varying
direction, the finite positive-definite inverse construction fails; it
should not be repaired by silently substituting an arbitrary inverse.
If breeding value has no variation at all, the variance share is undefined.

\section{Work through the scalar checkpoint}
\label{conversion:scalar_example}

Take the guide's values
\[
 P=1.5,\qquad \beta=0.4,\qquad \gamma=-0.2,\qquad G=0.5.
\]
Normalisation first gives
\[
 \alpha_w=1-\frac12\gamma P
 =1-\frac12(-0.2)(1.5)=1.15.
\]
The selected phenotype mean is \(P\beta=0.6\), its second moment is
\[
 P+P^2\gamma=1.5+(1.5)^2(-0.2)=1.05,
\]
and its variance is \(1.05-0.6^2=0.69\).

For breeding value, the conditional mean coefficient is
\(G/P=1/3\) and the conditional variance is
\[
 G-\frac{G^2}{P}
 =\frac12-\frac{1/4}{3/2}
 =\frac13.
\]
The selected mean and second moment are therefore
\begin{align*}
 m&=\frac GP(P\beta)=G\beta=0.2,\\
 \E_{\rm sel}[X^2]
 &=G-\frac{G^2}{P}+\frac{G^2}{P^2}(P+P^2\gamma)\\
 &=G+G^2\gamma=0.5+0.25(-0.2)=0.45.
\end{align*}
Subtracting the square of the mean gives
\[
 G_{\rm sel}=0.45-0.2^2=0.41.
\]
Only now invert the selected moments:
\begin{align}
 b&=\frac{0.2}{0.41}=\frac{20}{41}
     =0.487804878\ldots,\nonumber\\
 H&=\frac1{0.5}-\frac1{0.41}
     =-\frac{18}{41}=-0.439024390\ldots.
 \label{conversion:scalar_numbers}
\end{align}
Substituting them back gives
\[
 (G^{-1}-H)^{-1}
 =\left(2+\frac{18}{41}\right)^{-1}=0.41,
 \qquad G_{\rm sel}b=0.2.
\]
This is an algebraic check of the inversion. The fitted polynomial
\(1.15+0.4x-0.1x^2\) becomes negative sufficiently far into the tails.
The checkpoint is therefore a calculation with compatible regression
moments, not a proposal to use that polynomial as a nonnegative fitness
function over an entire Gaussian population.

\section{Why taking the logarithm of the regression is different}
\label{conversion:raw_log}

Let the fitted phenotype polynomial be
\[
 \widehat w(\vect x)
 =\alpha_w+\vect\beta^{\mathsf T}\vect x
       +\tfrac12\vect x^{\mathsf T}\bgamma\vect x.
\]
If \(\alpha_w>0\), continuity makes it positive in a neighbourhood of zero.
Within that neighbourhood,
\[
 \nabla\log\widehat w=\frac{\nabla\widehat w}{\widehat w}.
\]
Differentiate this product of \(\widehat w^{-1}\) and its gradient:
\[
 \nabla^2\log\widehat w
 =\frac{\nabla^2\widehat w}{\widehat w}
  -\frac{(\nabla\widehat w)(\nabla\widehat w)^{\mathsf T}}
         {\widehat w^2}.
\]
At zero, \(\widehat w=\alpha_w\),
\(\nabla\widehat w=\vect\beta\), and
\(\nabla^2\widehat w=\bgamma\). Thus
\begin{equation}
 \nabla\log\widehat w(\vect0)=\frac{\vect\beta}{\alpha_w},
 \qquad
 \nabla^2\log\widehat w(\vect0)
 =\frac{\bgamma}{\alpha_w}
   -\frac{\vect\beta\vect\beta^{\mathsf T}}{\alpha_w^2}.
 \label{conversion:raw_derivatives}
\end{equation}
In the scalar checkpoint these give
\[
 b_{\rm raw}=\frac{0.4}{1.15}=\frac8{23}=0.347826\ldots,
 \qquad
 H_{\rm raw}=-\frac{156}{529}=-0.294896\ldots.
\]
They differ from Equation~\eqref{conversion:scalar_numbers} because they
differentiate the fitted phenotype-space projection. They do not perform
environmental averaging or transfer selected moments to breeding value.
Positivity near the mean permits these local derivatives; it does not make
the polynomial positive over the whole phenotype distribution.

\section{Conditions to retain when using the result}

The conversion combines assumptions that belong to different steps. Gaussian
phenotype moments and regression orthogonality supply selected phenotype
moments. Joint Gaussian breeding values and phenotypes, the stated
cross-covariance, and selection mediation transfer them to breeding value.
Compatible centring, units and the half-quadratic coefficient convention
make the matrix operations refer to the same quantities. Positive selected
covariance permits inversion. Exact log-quadraticity is the additional
condition that turns matched parameters into the true local derivatives.

Uncertainty in \(\vect\beta,\bgamma,\mat G\), phenotype scaling and population
centring must be propagated together. Applying the conversion to joint
bootstrap or model draws preserves their dependencies; transforming separate
interval endpoints generally does not. Near-singular selected covariance can
make inverse-based parameters sensitive to small input changes. That
statistical uncertainty is separate from the validity of the underlying
moment and mediation assumptions.

The group \codepath{phenotype_regression_conversion} in
\codepath{code/reproduce_result_checks.py} checks the scalar
checkpoint and independently fitted positive-fitness examples. Its
multivariate calculation retains off-diagonal coefficients and noncommuting
matrix products. The calculation verifies those specified examples; the
derivations here state the conditions required to use the same route with
an existing selection regression.
